% GENERATED FILE. Edit canonical lessons, then run npm run generate:books.
% canonical-source-hash: fnv1a64:c9953f173c699101
% canonical-lessons: AR-C382-ihtijaj, AR-C382-karama, AR-C382-tazahara, AR-C382-ihtajja, AR-C382-idrab

\chapter{Protest}
\label{ch:ar-protest}
\hlchaptermodality{382}

\begin{chapteropening}
\textbf{By the end of this chapter:} \emph{I can talk about a protest, dignity and a strike, and say that people demonstrated or protested.}

Complete the last lesson of chapter 382: I can talk about a protest, dignity and a strike, and say that people demonstrated or protested.
\end{chapteropening}

\section[\texorpdfstring{\ar{احتجاج} (iḥtijāj)}{احتجاج (iḥtijāj)}]{\ar{احتجاج} (iḥtijāj) — a protest}

\label{lesson:AR-C382-ihtijaj}



\begin{quote}
Before the new one: say the Arabic for violence, then the Arabic for to suffer.
\end{quote}

\subsection*{You'll want to know: \ar{احتجاج}}
\textbf{\ar{احتجاج}} — \emph{iḥtijāj} — \textquotedblleft{}a protest\textquotedblright{}.

\begin{quote}
\ar{خرج} \ar{العمال} \ar{في} \ar{احتجاج} \ar{على} \ar{الرواتب}.
\end{quote}

\begin{quote}
\emph{kharaja al-ʿummāl fī iḥtijāj ʿalā ar-rawātib.}
\end{quote}

\begin{quote}
The workers came out in protest over wages.
\end{quote}

\begin{cousinweb}
\textbf{\ar{احتجاج}} (\emph{iḥtijāj}) is built on the native Arabic root \emph{ḥ-j-j}, \textquotedblleft{}to make for; to overcome in argument\textquotedblright{} (Lane). \emph{Iḥtijāj} is the Form VIII verbal noun, bringing an argument against: a protest.
\end{cousinweb}

\begin{grammarlens}[title={What you've built}]
One more word for giving reasons and opinions.
\end{grammarlens}

\subsection*{Your turn}
\begin{itemize}
\raggedright
  \item \emph{Say it:} \emph{iḥtijāj}
  \item \emph{Say it:} \emph{iḥtijāj}, once more
  \item \emph{Say it:} \emph{kharaja al-ʿummāl fī iḥtijāj ʿalā ar-rawātib.} — The workers came out in protest over wages.
  \item \emph{Recall:} say the Arabic for violence, then the Arabic for to suffer, then say \emph{iḥtijāj} again
\end{itemize}

\begin{tcolorbox}[breakable,colback=teal!4,colframe=teal!35!black,title={Before you move on}]
What does \ar{احتجاج} mean? (\textquotedblleft{}A protest\textquotedblright{}.) Say it once more.
\end{tcolorbox}
\section[\texorpdfstring{\ar{كرامة} (karāma)}{كرامة (karāma)}]{\ar{كرامة} (karāma) — dignity}

\label{lesson:AR-C382-karama}



\begin{quote}
Before the new one: say the Arabic for to suffer, then the Arabic for a protest.
\end{quote}

\subsection*{You'll want to know: \ar{كرامة}}
\textbf{\ar{كرامة}} — \emph{karāma} — \textquotedblleft{}dignity\textquotedblright{}.

\begin{quote}
\ar{لكل} \ar{إنسان} \ar{كرامة}.
\end{quote}

\begin{quote}
\emph{li-kulli insān karāma.}
\end{quote}

\begin{quote}
Every human being has dignity.
\end{quote}

\begin{cousinweb}
\textbf{\ar{كرامة}} (\emph{karāma}) is built on the native Arabic root \emph{k-r-m}, \textquotedblleft{}to be noble, generous\textquotedblright{} (Lane). \emph{Karāma} is nobility: dignity. \ar{كرم}, generosity, is from the same root.
\end{cousinweb}

\begin{grammarlens}[title={What you've built}]
One more word for giving reasons and opinions.
\end{grammarlens}

\subsection*{Your turn}
\begin{itemize}
\raggedright
  \item \emph{Say it:} \emph{karāma}
  \item \emph{Say it:} \emph{karāma}, once more
  \item \emph{Say it:} \emph{li-kulli insān karāma.} — Every human being has dignity.
  \item \emph{Recall:} say the Arabic for to suffer, then the Arabic for a protest, then say \emph{karāma} again
\end{itemize}

\begin{tcolorbox}[breakable,colback=teal!4,colframe=teal!35!black,title={Before you move on}]
What does \ar{كرامة} mean? (\textquotedblleft{}Dignity\textquotedblright{}.) Say it once more.
\end{tcolorbox}
\section[\texorpdfstring{\ar{تظاهر} (taẓāhara)}{تظاهر (taẓāhara)}]{\ar{تظاهر} (taẓāhara) — to demonstrate}

\label{lesson:AR-C382-tazahara}



\begin{quote}
Before the new one: say the Arabic for a protest, then the Arabic for dignity.
\end{quote}

\subsection*{You'll want to know: \ar{تظاهر}}
\textbf{\ar{تظاهر}} — \emph{taẓāhara} — \textquotedblleft{}to demonstrate\textquotedblright{}.

\begin{quote}
\ar{تظاهر} \ar{الطلاب} \ar{أمام} \ar{الجامعة}.
\end{quote}

\begin{quote}
\emph{taẓāhara aṭ-ṭullāb amāma al-jāmiʿa.}
\end{quote}

\begin{quote}
The students demonstrated in front of the university.
\end{quote}

Present: \ar{يتظاهر}, \emph{yataẓāhar}.

Also to pretend: \ar{تظاهر} \ar{بالنوم}, pretended to sleep.

\begin{cousinweb}
\textbf{\ar{تظاهر}} (\emph{taẓāhara}) is built on the native Arabic root \emph{ẓ-h-r}, \textquotedblleft{}to appear, to come out\textquotedblright{} (Lane). Lane gives Form VI for appearing and for people backing one another; modern Arabic uses it for a public show, to demonstrate, and with \ar{بـ} for a false show, to pretend.
\end{cousinweb}

\begin{grammarlens}[title={What you've built}]
One more word for giving reasons and opinions.
\end{grammarlens}

\subsection*{Your turn}
\begin{itemize}
\raggedright
  \item \emph{Say it:} \emph{taẓāhara}
  \item \emph{Say it:} \emph{taẓāhara}, once more
  \item \emph{Say it:} \emph{taẓāhara aṭ-ṭullāb amāma al-jāmiʿa.} — The students demonstrated in front of the university.
  \item \emph{Recall:} say the Arabic for a protest, then the Arabic for dignity, then say \emph{taẓāhara} again
\end{itemize}

\begin{tcolorbox}[breakable,colback=teal!4,colframe=teal!35!black,title={Before you move on}]
What does \ar{تظاهر} mean? (\textquotedblleft{}To demonstrate\textquotedblright{}.) Say it once more.
\end{tcolorbox}
\section[\texorpdfstring{\ar{احتج} (iḥtajja)}{احتج (iḥtajja)}]{\ar{احتج} (iḥtajja) — to protest}

\label{lesson:AR-C382-ihtajja}



\begin{quote}
Before the new one: say the Arabic for dignity, then the Arabic for to demonstrate.
\end{quote}

\subsection*{You'll want to know: \ar{احتج}}
\textbf{\ar{احتج}} — \emph{iḥtajja} — \textquotedblleft{}to protest\textquotedblright{}.

\begin{quote}
\ar{احتج} \ar{السكان} \ar{على} \ar{إغلاق} \ar{الحديقة}.
\end{quote}

\begin{quote}
\emph{iḥtajja as-sukkān ʿalā ighlāq al-ḥadīqa.}
\end{quote}

\begin{quote}
The residents protested against closing the park.
\end{quote}

Present: \ar{يحتج}, \emph{yaḥtajj}.

It takes \ar{على}: \ar{احتج} \ar{على} ….

\begin{cousinweb}
\textbf{\ar{احتج}} (\emph{iḥtajja}) is built on the native Arabic root \emph{ḥ-j-j}, \textquotedblleft{}to make for; to overcome in argument\textquotedblright{} (Lane). It is Form VIII, to bring an argument against: to protest.
\end{cousinweb}

\begin{grammarlens}[title={What you've built}]
One more word for giving reasons and opinions.
\end{grammarlens}

\subsection*{Your turn}
\begin{itemize}
\raggedright
  \item \emph{Say it:} \emph{iḥtajja}
  \item \emph{Say it:} \emph{iḥtajja}, once more
  \item \emph{Say it:} \emph{iḥtajja as-sukkān ʿalā ighlāq al-ḥadīqa.} — The residents protested against closing the park.
  \item \emph{Recall:} say the Arabic for dignity, then the Arabic for to demonstrate, then say \emph{iḥtajja} again
\end{itemize}

\begin{tcolorbox}[breakable,colback=teal!4,colframe=teal!35!black,title={Before you move on}]
What does \ar{احتج} mean? (\textquotedblleft{}To protest\textquotedblright{}.) Say it once more.
\end{tcolorbox}
\section[\texorpdfstring{\ar{إضراب} (iḍrāb)}{إضراب (iḍrāb)}]{\ar{إضراب} (iḍrāb) — a strike}

\label{lesson:AR-C382-idrab}



\begin{quote}
Before the new one: say the Arabic for to demonstrate, then the Arabic for to protest.
\end{quote}

\subsection*{You'll want to know: \ar{إضراب}}
\textbf{\ar{إضراب}} — \emph{iḍrāb} — \textquotedblleft{}a strike\textquotedblright{}.

\begin{quote}
\ar{أعلن} \ar{العمال} \ar{إضرابا}.
\end{quote}

\begin{quote}
\emph{aʿlana al-ʿummāl iḍrāban.}
\end{quote}

\begin{quote}
The workers announced a strike.
\end{quote}

\begin{cousinweb}
\textbf{\ar{إضراب}} (\emph{iḍrāb}) is built on the native Arabic root \emph{ḍ-r-b}, \textquotedblleft{}to strike, to beat\textquotedblright{} (Lane). \emph{Iḍrāb} is the Form IV verbal noun; \emph{aḍraba ʿan} is to turn away from a thing, and so to refuse to work: a strike.
\end{cousinweb}

\begin{grammarlens}[title={What you've built}]
One more word for giving reasons and opinions.
\end{grammarlens}

\subsection*{Your turn}
\begin{itemize}
\raggedright
  \item \emph{Say it:} \emph{iḍrāb}
  \item \emph{Say it:} \emph{iḍrāb}, once more
  \item \emph{Say it:} \emph{aʿlana al-ʿummāl iḍrāban.} — The workers announced a strike.
  \item \emph{Recall:} say the Arabic for to demonstrate, then the Arabic for to protest, then say \emph{iḍrāb} again
\end{itemize}

\begin{tcolorbox}[breakable,colback=teal!4,colframe=teal!35!black,title={Before you move on}]
What does \ar{إضراب} mean? (\textquotedblleft{}A strike\textquotedblright{}.) Say it once more.
\end{tcolorbox}
